%!TEX encoding = UTF-8 Unicode
%!TEX root = ./../main.tex
%!TEX TS-program = xelatex

\begin{center}
\begin{figure}[!h]
  	\centering
 	\includegraphics[width=1.5cm]{../pictures/frontespizio/logo_unipa_piccolo.png}
\end{figure}
\textsc{\textbf{UNIVERSIT\`A DEGLI STUDI DI PALERMO}} \\
SCUOLA POLITECNICA\\
\small{Laurea Magistrale in Ingegneria Informatica}\\
\end{center}
\begin{center}
\uppercase{Estrazione di conoscenza mediante algoritmi di apprendimento attivo per l'inferenza regolare}
%SPERIMENTAZIONE DI ALGORITMI DI ACTIVE LEARNING NELL'INFERENZA INDUTTIVA REGOLARE
\end{center}

\begin{minipage}[t]{0.65\textwidth}\raggedright
{ \scriptsize{TESI DI LAUREA DI:}\\
\scriptsize{\textbf{Dott. Nicola Ciaco}}
}
\end{minipage}
\begin{minipage}[t]{0.47\textwidth}\raggedright
{\scriptsize{RELATORE:}\\
\vspace{0.7mm}
\scriptsize{\textbf{Chiar.mo Prof. Salvatore Gaglio}}\\
\vspace{0.7mm}
\scriptsize{CORRELATORI:}\\
\vspace{0.7mm}
\scriptsize{\textbf{Ing. M. Ortolani, Ing. P. Cottone}}\\
%\scriptsize{\textbf{Ing. Pietro Cottone}}
%\vspace{3mm}
}
\end{minipage}
\begin{center}
{\small Anno Accademico 2017/18 \\}%inserire l'anno accademico a cui si è iscritti
\vspace{3mm}
\small{\textbf{Sommario}}\\
\end{center}


\label{cap:sommario}
Un sistema intelligente è caratterizzato dalla capacità di apprendere in modo automatico, la quale a sua volta presuppone la capacità sia di rappresentare la conoscenza nota a priori sia di \textit{inferirne} di nuova. Allo stato dell'arte attuale, le tecniche più efficaci per l'\textit{estrazione di conoscenza} sono state sviluppate nell'ambito di quella che è nota come Statistical Learning Theory (SLT). Per quanto tali metodi forniscano risultati apprezzabili in termini di accuratezza, essi presentano tuttavia anche dei non trascurabili svantaggi, quali ad esempio il fatto che il loro funzionamento si basa sull'inferire dei parametri ottimali per un modello che è visto necessariamente come una \textit{black-box}, rendendo ardua o impossibilitando di fatto l'\textit{interpretazione} dei campioni a partire dai quali il modello è stato costruito.

La presente tesi si colloca invece nel solco di un framework teorico alternativo, noto come Algorithmic Learning Theory (ALT), secondo cui i dati non sono considerati come campioni casuali da mappare in uno spazio vettoriale, bensì come specifiche istanze del modello nascosto che è oggetto di inferenza. In particolare, lo studio si è concentrato su una tecnica di estrazione di conoscenza che va sotto il nome di \textit{Inferenza Grammaticale} (GI), un processo di apprendimento che si basa sull'\textit{induzione} nel contesto dei linguaggi formali e del relativo formalismo grammaticale. In tale contesto, l'obiettivo è selezionare il migliore modello rappresentato mediante il tipico riconoscitore di una grammatica formale, ovvero l'automa a stati finiti, consistente con i campioni iniziali, che a loro volta sono interpretabili come stringhe generabili dalla grammatica. Nella fattispecie, la classe dei linguaggi oggetto di studio per la presente tesi è quella dei linguaggi regolari e si parla pertanto più specificamente di \textit{Inferenza Induttiva Regolare} (IIR), i cui limiti teorici sono stati messi in luce dal lavoro di Gold \cite{Gold67}. Gli algoritmi della letteratura iniziale sull'argomento riconducevano l'apprendimento ad una ricerca euristica in un spazio rappresentato come un grafo contenente gli automi consistenti con i campioni forniti in cui il modello, inizialmente iperspecializzato, viene progressivamente generalizzato. Tale approccio va sotto il nome di \textit{passive learning} ed è caratterizzato da un'inevitabile esplosione combinatoria che ne rende la complessità ingestibile a meno di non sacrificare le garanzie teoriche di terminazione e ottimalità.

Nella presente tesi si è quindi scelto di seguire un approccio proposto nella letteratura più recente, duale rispetto al precedente, secondo cui un modello, inizialmente molto generale e quindi poco accurato, viene progressivamente specializzato per rappresentare con precisione i dati forniti. Questo paradigma, noto come \textit{active learning}, presuppone l'esistenza di un \textit{oracolo} che guida l'apprendimento rispondendo ad alcuni tipi di query sottopostegli attivamente dal sistema che cerca di inferire il modello. La base teorica del paradigma è stata fornita dagli studi di Angluin \cite{Angluin87} che stabiliscono che l'oracolo debba appartenere alla classe dei cosiddetti Minimally Adequate Teacher (MAT) che, per garantire l'abilità di fornire risposte utili ai principali tipi di query considerate, richiede la conoscenza preliminare dell'automa oggetto d'inferenza.
Questo è un requisito molto forte e fin troppo stringente, che limita l'utilizzo in contesti reali dell'inferenza induttiva regolare declinata nell'accezione dell'\textit{active learning}.
 

Il cuore del lavoro svolto in questa tesi ha quindi riguardato la sostituzione di un tradizionale oracolo con una sua approssimazione mediante un classificatore statistico costruito a partire da esempi positivi e negativi del linguaggio target nell'ottica di permetterne un utilizzo nelle applicazioni reali. L'obiettivo è quello di tracciare un parallelo tra i modelli ottenibili seguendo il paradigma della teoria dell'apprendimento statistico e quelli formulabili algoritmicamente e di verificare sperimentalmente se l'uso combinato possa condurre al superamento dei limiti di entrambi. In altre parole, si vuole costruire per i campioni dati un modello strutturale, da usare in luogo dell'oracolo in modo da superarne i requisiti stringenti.  Il modello statistico prescelto per affrontare il delineato problema di classificazione binaria è stato quello delle Support Vector Machines (SVM) perchè rappresentano un modello molto potente. Inoltre malgrado siano state proposte in letteratura delle loro applicazioni nel contesto dei linguaggi formali esse non sono state contestualizzate nell’ambito dell’inferenza induttiva regolare mediante \textit{active learning} oppure erano limitate all'apprendimento di specifiche classi di linguaggi sottoinsiemi di quelli regolari.


Al fine di valutare la correttezza della soluzione ottenuta, durante il lavoro di tesi si è progettata un'implementazione in linguaggio C++11, integrando il codice nella preesistente libreria di passive learning Gi-learning \cite{Cot16}. L'implementazione è stata testata sperimentalmente inizialmente su una nota classe di automi che rappresentano linguaggi semplici, ovvero i Tomita \cite{Tomita82, Dupont94}, ottenendo ottimi risultati paragonabili al caso ideale.
Infine le performance e la precisione dei modelli ottenuti dall'approccio qui proposto sono state vagliate su data set di automi estratti casualmente, la cui complessità è paragonabile ai sistemi da apprendere nei casi pratici. I risultati ottenuti rivelano che all'aumentare della complessità del target si ha un progressivo decadimento del grado di approssimazione ed è quindi possibile concludere che un livello di fiducia alto sull'utilizzo dell'algoritmo è ottenibile solo fino ad un certo livello di complessità del target.


